\documentclass{beamer}
%[aspectratio=169]   \usepackage[czech]{babel}
\usepackage{apo-lecture-cz}
\usepackage{pdfpages}
\usepackage{pdfcomment}
\usepackage{listings}
\usepackage{array,multirow}

\subtitle{Lekce 04. Hierarchie paměti}
\author{Pavel Píša \phantom{xxxxxxx} Petr Štěpán \\ \small\texttt{pisa@fel.cvut.cz}\phantom{xxxx}\small\texttt{stepan@fel.cvut.cz}}
\begin{document}

\maketitle

\section{Paměť úvod}

\input{common/memory/tex/memory-introduction-cz.tex}

\section{Technologie polovodičové paměťi}

\input{common/memory/tex/memory-kinds-cz.tex}

\input{common/memory/tex/memory-dimm-sdram-cz.tex}

\input{common/memory/tex/memory-gap-cz.tex}

\section{Zrychlení přístupu k datům, hierarchie, vyrovnávací paměť}

\input{common/memory/tex/memory-hierarchy-cz.tex}

\input{common/memory/tex/memory-cache-basis-cz.tex}

\input{common/simulator/tex/qtrvsim-cache-examples-cz.tex}

\input{common/memory/tex/memory-cache-parameter-impact-cz.tex}

\section{Virtuální paměť a stránkování}

\input{common/memory/tex/memory-paging-cz.tex}

\section{Cache a stránkování dohromady}

\input{common/memory/tex/memory-64bit-cache-and-paging-cz.tex}

\begin{frame}
\frametitle{Odkazy a literatura}

\begin{itemize}
\item Ulrich Drepper: \href{https://lwn.net/Articles/250967/}{What every programmer should know about memory}
\item Agner Fog: \href{https://www.agner.org/optimize/}{Software optimization resources. C++ and assembly}
\item \url{https://www.7-cpu.com/cpu/Haswell.html}
\item \url{https://www.7-cpu.com/cpu/Skylake.html}
\item WikiChips: \href{https://en.wikichip.org/wiki/amd/microarchitectures/zen}{Zen - Microarchitectures - AMD}

\end{itemize}

\end{frame}

\begin{frame}
\frametitle{Prostor na poznámky}

\end{frame}

\end{document}

